
        \documentclass[fleqn]{article}      
        \usepackage{amsmath}
        \usepackage{fontspec}
        \usepackage{kotex} % 한국어 지원  

        \begin{document}      
        \textbf{쉬운 문제}  
\noindent 1. 다음 행렬 \( A = \begin{pmatrix} 1 & 2 \\ 3 & 4 \end{pmatrix} \)의 행렬식은 무엇인가?  
\\\\  
\noindent 2. 행렬 \( B = \begin{pmatrix} 5 & 6 \\ 7 & 8 \end{pmatrix} \)의 전치행렬 \( B^T \)는 무엇인가?  
\\\\  
\noindent 3. 두 행렬 \( C = \begin{pmatrix} 1 & 0 \\ 0 & 1 \end{pmatrix} \)과 \( D = \begin{pmatrix} 2 & 2 \\ 2 & 2 \end{pmatrix} \)의 합은?  
\\\\  
\noindent 4. 다음 행렬 \( E = \begin{pmatrix} 3 & 5 \\ 1 & 2 \end{pmatrix} \)의 (1, 2)번째 원소는 무엇인가?  
\\\\  
\noindent 5. 행렬 \( F = \begin{pmatrix} 1 & 2 \\ 3 & 4 \end{pmatrix} \)의 첫 번째 행의 합은 얼마인가?  
\\\\  

\textbf{중간 문제}  
\noindent 6. 행렬 \( G = \begin{pmatrix} 1 & 4 \\ 2 & 3 \end{pmatrix} \)의 고유값을 구하라.  
\\\\  
\noindent 7. 행렬 \( H = \begin{pmatrix} 2 & 1 \\ 5 & 3 \end{pmatrix} \)에 대해 \( H^{-1} \)를 구하라.  
\\\\  
\noindent 8. 두 행렬 \( J = \begin{pmatrix} 1 & 3 \\ 2 & 4 \end{pmatrix} \)과 \( K = \begin{pmatrix} 5 & 6 \\ 7 & 8 \end{pmatrix} \)의 곱을 구하라.  
\\\\  
\noindent 9. 행렬 \( L = \begin{pmatrix} 2 & 3 \\ 5 & 7 \end{pmatrix} \)의 행렬식을 계산하라.  
\\\\  
\noindent 10. 행렬 \( M = \begin{pmatrix} 1 & 2 \\ 3 & 4 \end{pmatrix} \)와 스칼라 3의 곱은?  
\\\\  

\textbf{어려운 문제}  
\noindent 11. \( N = \begin{pmatrix} 2 & 0 & 1 \\ 1 & 3 & 0 \\ 0 & 2 & 1 \end{pmatrix} \)의 고유값을 구하라.  
\\\\  
\noindent 12. 두 행렬 \( P = \begin{pmatrix} 1 & 2 \\ 3 & 4 \end{pmatrix} \)과 \( Q = \begin{pmatrix} 0 & 1 \\ 1 & 0 \end{pmatrix} \)의 곱 \( PQ \)와 \( QP \)를 구하고, 두 결과가 같은지 확인하라.  
\\\\  
\noindent 13. \( R = \begin{pmatrix} 4 & 2 \\ 2 & 3 \end{pmatrix} \)의 고유 벡터를 구하라.  
\\\\  
\noindent 14. \( S = \begin{pmatrix} 1 & 1 & 1 \\ 0 & 1 & 2 \\ 1 & 0 & 1 \end{pmatrix} \)에서 \( S \)의 역행렬을 구하라.  
\\\\  
\noindent 15. 행렬 \( T = \begin{pmatrix} 0 & 1 \\ -1 & 0 \end{pmatrix} \)의 제곱 \( T^2 \)를 구하라.  
\\\\  

\textbf{답안}  
1. -2 \\\\\\\  
2. \( \begin{pmatrix} 5 & 7 \\ 6 & 8 \end{pmatrix} \) \\\\\\\  
3. \( \begin{pmatrix} 3 & 2 \\ 2 & 2 \end{pmatrix} \) \\\\\\\  
4. 5 \\\\\\\  
5. 3 \\\\\\\  
6. 5, -1 \\\\\\\  
7. \( \begin{pmatrix} 3 & -1 \\ -5 & 2 \end{pmatrix} \) \\\\\\\  
8. \( \begin{pmatrix} 26 & 36 \\ 34 & 46 \end{pmatrix} \) \\\\\\\  
9. -1 \\\\\\\  
10. \( \begin{pmatrix} 3 & 6 \\ 9 & 12 \end{pmatrix} \) \\\\\\\  
11. 3, 2, 1 \\\\\\\  
12. \( PQ = \begin{pmatrix} 2 & 4 \\ 3 & 4 \end{pmatrix} \), \( QP = \begin{pmatrix} 3 & 2 \\ 4 & 1 \end{pmatrix} \) (다르다) \\\\\\\  
13. \( \begin{pmatrix} 1 \\ -1 \\ 0 \end{pmatrix} \), \( \begin{pmatrix} 2 \\ 1 \\ 0 \end{pmatrix} \) \\\\\\\  
14. \( \begin{pmatrix} 1 & -1 & 1 \\ 0 & 1 & -2 \\ -1 & 0 & 1 \end{pmatrix} \) \\\\\\\  
15. \( \begin{pmatrix} -1 & 0 \\ 0 & -1 \end{pmatrix} \) \\\\\\\        
        \end{document} 
    